\paragraph{R4-O4, R4-W3 (Multimodal appears closer to shared-embedding retrieval; explain cross-modal constraints and consistency validation).}

\begin{itemize}[leftmargin=1.2em, itemsep=1pt, topsep=2pt]
\item \textit{O4: ``\dots it appears closer to shared-embedding-based multimodal retrieval rather than true multimodal counterfactual reasoning. The authors are encouraged to further explain how cross-modal constraints and consistency validation are handled.''}
\item \textit{W3 (part): ``The multimodal contribution \dots would benefit from more detailed justification and discussion.''}
\end{itemize}
\smallskip

\noindent\textbf{Response.}
We agree and state the scope honestly: the current prototype uses a shared embedding (CLIP~\cite{radford2021learning}) as a proposal mechanism with post-filtering, consistent with Section~3, which already holds that similarity search \emph{``serves only as a proposal mechanism, while correctness is determined through cross-modal validation that enforces feasibility and consistency,''} making multimodal counterfactual querying \emph{``inherently verification-driven.''} We make the mechanism concrete: a per-modality constraint table (Table~\ref{tab:modality}) and a cross-modal consistency predicate within $\omega$~\cite{chen2019cross}, whereby a candidate must satisfy each modality's constraints and be jointly consistent across views. We elevate true multimodal reasoning (constraints integrated into retrieval, not post-filtering) to an explicit research challenge in the success criteria. The important distinction is that DcCA exposes modality constraints as declarative query semantics: the current prototype realizes them through validation, but the abstraction enables future optimizers to push these constraints into indexing and retrieval, analogous to how relational databases evolved from filtering after scans to predicate-aware query optimization.

\noindent\emph{Revision to Section~3 (cross-modal consistency):}
\begin{quote}
\textcolor{blue}{For multimodal queries, $\omega$ includes a cross-modal consistency predicate: a candidate must satisfy each modality's admissibility and be jointly consistent across modalities (e.g., a tabular edit must imply a consistent text/image change). Shared-embedding retrieval proposes candidates; cross-modal validation confirms consistency, so correctness is verification-driven. Integrating these constraints into index construction, rather than applying them as a post-filter, is a key open challenge.}
\end{quote}
